% !TEX root = ../main.tex
\chapter{문헌 검토 및 이론적 틀}
\label{ch:literature}

본 장은 그래프 기반 온체인 평판 점수 산출, ERC-20 allowance 의미론, GMX V2 거래 성공 검증을 검토하고, \ref{ch:methodology}장과 \ref{ch:results}장의 주요 EndorseRank 대 AWP 비교를 동기화하는 이론적 틀---하이퍼링크 전파, 도메인 정렬, 계산 복잡도---을 전개한다.

\dissertationSubheading[sec:graph-reputation]{그래프 기반 온체인 평판}

PageRank 및 변형은 방향성 엣지를 통해 전파되는 영향력으로 노드를 순위 매긴다(Page et al., 1998). Do et al.\ (2023)은 Ethereum 거래 그래프에 Adaptive Weighted PageRank(AWP)를 적용하여, 엣지를 송금액과 로지스틱 시간 감쇠로 가중한다. 그들은 \textbf{in-degree}와 \textbf{in-value}---유입 송금의 횟수와 합---를 경제적 중요성에 대한 도메인 직관 프록시로 사용하여 점수를 검증한다. Lin et al.\ (2024)은 거래 연결성만으로는 리스크 관련 구조를 놓칠 수 있으며, 리스크 신호가 네트워크를 통해 전파되어야 함을 보여 주어, 순수 PageRank를 넘어선 정교한 그래프 의미론의 필요성을 뒷받침한다. DeFi 신용 점수 문헌에는 전망 이론 및 블록체인 기반 모델(Hassija et al., 2020)과 대출 프로토콜 머신러닝 파이프라인(Kotb, 2023)이 포함된다; Packin and Lev-Aretz (2024)는 가명성 하 감사 가능성 제약을 강조한다. EndorseRank는 블랙박스 예측 점수가 아닌 해석 가능한 그래프 전파 범위에 머문다.

\dissertationSubheading[sec:allowance-endorsement]{보증으로서의 ERC-20 allowance}

ERC-20 \texttt{approve}와 EIP-2612 \texttt{permit}은 명시적 지출 승인을 기록한다(Ethereum Improvement Proposals, 2015, 2020). EndorseRank는 소유자--지출자 쌍당 최신 allowance를 방향성 보증 그래프의 엣지 강도로 취급하고, 해당 스냅샷에 PageRank를 실행한다. 새 승인이 이전 승인을 덮어쓰므로 과거 시간 감쇠가 불필요하다.

\dissertationSubheading[sec:gmx-validation]{거래 성공 검증 프레임 (GMX V2)}

Arbitrum One에서 실현 PnL(\textbf{basePnlUsd})을 갖는 밀도 높은 지갑 귀속 \textbf{PositionDecrease} 이벤트는 거래 성공 프록시를 제공한다: 청산 성공 횟수, 실현 이익 대리변수, 청산 성공률(GMX, n.d.; Sch\"{a}r, 2021). 담보 부족 대출 부도 레이블은 본 체인에서 희소하다; GMX V2 무기한 청산은 도메인 직관 검증을 위한 밀도 높은 지갑 수준 실현 결과를 제공한다. 이 프록시들은 Do et al.의 in-degree/in-value 논리를 따르되, 순수 유입이 아닌 담보화 무기한 거래를 대상으로 한다.

\dissertationSubheading[sub:gf-pr]{보조 확장: 결과-네이티브 참조 (GF-PR)}

실증 단계에서 \textbf{GF-PR} (GainFlow PageRank)이 거래 결과 프록시 정렬 상한을 설정하는 보조 진단으로 추가되었다. GF-PR은 동일 GMX \texttt{PositionDecrease} 스트림에서 방향성 스타 그래프를 구축한다: 합성 이익 풀 노드가 양의 실현 PnL 지갑으로 엣지를 보내고, 지갑은 손실 싱크로 엣지를 보낸다(일반 손실 $|\text{PnL}|+0.1$, 청산 $|\text{PnL}|+1.0$). GF-PR에는 지갑-대-지갑 사회적 엣지가 없으며, EndorseRank나 AWP의 동등한 방법이 아니라 \emph{결과-네이티브 상한 참조}(SRQ1)로 기능한다. 거래 성공 프록시와의 높은 정렬은 구성상 예상되며, 우수한 사회적 평판이 아닌 구성개념 중첩으로 해석해야 한다.

\dissertationSubheading[sec:theoretical-framework]{이론적 틀}

본 연구는 하이퍼링크-인용 전파, 구성개념(도메인) 정렬, 계산 복잡도에 기반하여 EndorseRank와 AWP를 평판 구조(H1), 거래 성공 정렬(H2), 계산 효율(H3) 측면에서 비교한다.

\textbf{하이퍼링크--인용 의미론.} 원래 PageRank 알고리즘(Page et al., 1998)은 유입 하이퍼링크를 인용으로 취급한다. Do et al.\ (2023)은 시간 감쇠 엣지 가중치로 Ethereum 송금에 이 논리를 적용한다. EndorseRank는 allowance가 명시적 온체인 지출 권한 위임이기 때문에(Ethereum Improvement Proposals, 2015, 2020) ERC-20 approve/permit 엣지로 인용 기반 전파를 확장한다. 송금과 달리, 송금은 신뢰 의도 없이 거래나 라우팅을 반영할 수 있지만, allowance는 취소될 때까지 지속적인 승인을 기록한다. 그래프 기반 리스크 모델(Lin et al., 2024)은 네트워크 구조가 원시 활동 횟수를 넘어 정보를 담고 있음을 추가로 보여 주어, 송금량만이 아닌 방향성 보증 그래프를 지지한다.

\textbf{도메인 정렬 프레임워크.} 측정 이론에서 구성개념 정렬은 지표가 관찰 가능한 결과를 얼마나 잘 추적하는지 평가한다(Cronbach \& Meehl, 1955). 순위 기반 평판 점수에 대해 Spearman's $\rho$와 Kendall's $\tau$는 분류 임계값을 가정하지 않고 단조적·쌍별 순서 일치를 평가한다(Do et al., 2023). GMX V2의 청산 성공 횟수와 실현 이익 대리변수는 Arbitrum One에서 관찰 가능한 거래 성공 측정치로, 외부 대출 상환 근거를 주장하지 않으면서 도메인 직관 요건을 충족한다.

\textbf{계산 복잡도.} PageRank의 반복당 비용은 $O(|E|)$이다(Page et al., 1998). approve/permit 그래프는 송금 그래프보다 희소하다(엣지 $|E|$가 적음), 반복당 비용 감소, 더 빠른 수렴, 메모리 절감을 가져온다---H3의 근거이다.

그림~\ref{fig:conceptual-model}은 주요 EndorseRank--AWP 비교를 H1--H3 및 RQ1--RQ4에 매핑한 개념 모형을 제시한다. GF-PR(점선)은 실증 연구 중 추가된 SRQ1용 보조 결과-네이티브 상한이다.

\begin{figure}[htbp]
\centering
% !TEX root = ../main.tex
\begin{tikzpicture}[
  font=\small,
  box/.style={draw, rounded corners, align=center, minimum width=10.5cm, minimum height=1.0cm},
  smallbox/.style={draw, rounded corners, align=center, minimum width=3.4cm, minimum height=1.05cm},
  refbox/.style={draw, dashed, rounded corners, align=center, minimum width=3.2cm, minimum height=0.9cm},
  arrow/.style={->, thick}
]
\node[box] (center) at (0,2.8) {주요 비교: EndorseRank vs AWP};

\node[smallbox] (rep) at (-5,0.8) {평판\\구조};
\node[smallbox] (risk) at (0,0.8) {거래 성공\\정렬};
\node[smallbox] (eff) at (5,0.8) {계산\\효율};

\node[smallbox] (h1) at (-5,-1.2) {H1 / RQ1\\EndorseRank--AWP\\발산};
\node[smallbox] (h2) at (0,-1.2) {H2 / RQ2,4\\사회적 방법\\정렬};
\node[smallbox] (h3) at (5,-1.2) {H3 / RQ3\\EndorseRank\\효율};

\node[refbox] (gfref) at (0,-2.6) {보조: GF-PR\\결과-네이티브 상한 (SRQ1)};

\draw[arrow] (center.south) -- (rep.north);
\draw[arrow] (center.south) -- (risk.north);
\draw[arrow] (center.south) -- (eff.north);

\draw[arrow] (rep.south) -- (h1.north);
\draw[arrow] (risk.south) -- (h2.north);
\draw[arrow] (eff.south) -- (h3.north);

\draw[arrow, dashed] (h2.south) -- (gfref.north);
\end{tikzpicture}

\caption{개념 모형: 평판 구조(H1/RQ1), 사회적 방법 정렬(H2/RQ2--4), 계산 효율(H3/RQ3)에 매핑된 주요 EndorseRank vs.\ AWP 비교. GF-PR(점선)은 보조 결과-네이티브 상한 참조(SRQ1)이며, 세 번째 동등 사회적 방법이 아니다.}
\label{fig:conceptual-model}
\end{figure}

\textbf{GF-PR의 보조 역할.} H1과 H3는 EndorseRank와 AWP만을 다룬다. H2는 다섯 프록시 계열에서 EndorseRank와 AWP 간 사회적 방법 정렬을 비교한다(RQ2, RQ4). GF-PR은 결과 흐름에 직접 순위를 매길 때 달성 가능한 거래 결과 프록시 정렬의 최대치를 한정하여, 주요 가설 검정에 들어가지 않고 SRQ1에 대한 구성개념 중첩을 명확히 한다.

\dissertationSubheading[sub:terminology]{용어 안내서}

\ref{ch:introduction}장의 정의는 공식 구성개념을 확립한다; 다음 안내는 결과 해석과 흔한 오독을 피하는 데 도움을 준다:

\begin{itemize}
\item \textbf{Primary methods:} EndorseRank(allowance 그래프)와 AWP(송금 그래프). 둘 다 지갑-대-지갑 엣지의 사회적 PageRank 방법이며, 배포 가능한 평판 신호를 목표로 한다.
\item \textbf{Supplementary reference:} 합성 POOL/SINK 노드로/from GMX PnL 흐름의 GF-PR; SRQ1용 결과-네이티브 상한이며, 배포 가능한 사회적 점수가 아니다.
\item \textbf{Primary graph construct:} 방법이 순위를 매기도록 \emph{설계된} 그래프(예: EndorseRank의 $G_{\text{approve}}$). 점수 산출 후 적용되는 검증 프록시와 구별된다.
\item \textbf{Validation proxy:} 순위가 외부 구성개념과 정렬되는지 검증하는 관찰 가능 지갑 지표(예: in-degree, GMX 청산 성공 횟수)(Cronbach \& Meehl, 1955).
\item \textbf{Construct overlap:} 방법의 입력 그래프와 검증 프록시가 동일 데이터 출처를 공유할 때(GF-PR vs.\ 거래 결과 프록시), 사회적 평판을 함의하지 않고 $\tau$를 부풀린다.
\item \textbf{Alignment bands (Kendall $\tau$):} $|\tau|<0.10$ 무시 가능한 노이즈; $0.10$--$0.35$ 약~중간; $\geq 0.35$ 유사 PageRank/블록체인 연구에서 실질적으로 의미 있음; 값은 동일 코호트에서 방법 간 \emph{상대적} 비교에 사용된다.
\end{itemize}

\dissertationSubheading[sec:domain-alignment-metrics]{도메인 정렬 지표}

Do et al.\ (2023)과 Cronbach \& Meehl (1955)에 따라, 순위 기반 점수의 구성개념 정렬은 평판 점수와 프록시 순위 간 Spearman's $\rho$와 Kendall's $\tau$를 사용한다. 본 학위논문은 \textbf{다섯 프록시 계열}(송금, allowance, 역리스크, Sybil 조정 안정성, 거래 성공)을 평가한다; 운영 정의는 \ref{ch:methodology}장을 참조한다.
